\documentclass[10pt,a4paper]{amsart}
\usepackage[margin=19mm]{geometry}
\usepackage{amsmath,amssymb,mathtools}
\usepackage[scr=rsfs,cal=euler]{mathalfa}
\usepackage{xfrac}
\usepackage[unicode,hidelinks,bookmarks=false]{hyperref}
\newcommand{\ie}{i.e.\ }
\newcommand{\cf}{cf.\ }
\newcommand{\resp}{respectively\ }
\newcommand{\Subsection}{\subsection}
\providecommand{\nobreakdash}{\nobreak\hbox{-}}
\setlength{\emergencystretch}{2em}
\multlinegap0pt
\usepackage{bm}
\usepackage{bbm}
\usepackage{dsfont}
\usepackage{xspace}
\usepackage{cancel}
\usepackage{mathtools}
\usepackage{amsthm}
\makeatletter
\@ifundefined{claim}{\newtheorem{claim}{Claim}}{}
\@ifundefined{lemma}{\newtheorem{lemma}{Lemma}}{}
\@ifundefined{theorem}{\newtheorem{theorem}{Theorem}}{}
\makeatother
\newcommand\FUT{\mathsf{FUT}}
\newcommand\Nat{\mathbb{N}}
\newcommand\RCA{\mathsf{RCA}}
\newcommand\WOP{\mathsf{WOP}}
\newcommand\WO{\mathsf{WO}}
\newcommand\W{\mathrm{W}}
\newcommand\boldomega{\pmb{\omega}}
\newcommand{\LoX}{\mathcal{X}}
\title[Reductions of well-ordering principles to combinatorial theo]{Reductions of well-ordering principles to combinatorial theorems}
\author{Lorenzo Carlucci, Leonardo Mainardi, Konrad Zdanowski}
\date{Source: arXiv:2401.04451v1 (CC BY 4.0)}
\begin{document}
\maketitle

\noindent\emph{Source and licence.} Reductions of well-ordering principles to combinatorial theorems. Version arXiv:2401.04451v1 (CC BY 4.0), distributed under the Creative Commons Attribution 4.0 International licence, as recorded in the arXiv licence metadata for this exact version and shown on the abstract page https://arxiv.org/abs/2401.04451v1. Full rights notice: licenses/arxiv-2401.04451-CC-BY-4.0.txt.\par\smallskip

\noindent\textit{Editorial scope.} Counted unit: the Theorem whose source label is \texttt{prop:RTomega} in the article (source0.tex lines 883--886, source label \texttt{prop:RTomega}), delivered together with the article's own complete original proof of it (source0.tex lines 889--972, one \texttt{proof} environment of 84 lines). No mathematics is written, repaired or re-derived: statement and proof are the article's own text line for line. Required context retained uncounted: lem:decribile-terms (lines 855--857). Every definition, display and label of those lines is retained unchanged. Omitted from this package and not redistributed: 1--854, 858--882, 887--888, 973--1113. Nothing inside a retained excerpt is omitted or altered: every retained body is delivered exactly as the article's lines are written, so no omission receipt of any kind accompanies this package. Citations: the retained excerpts carry no citation command at all, so no bibliography entry of the article is transcribed and no reference is redistributed. Reference adaptation: none was needed and none was made; every reference inside the delivered unit resolves inside it, so no unresolved reference or citation is produced. Printed slips kept literal: no printed word, symbol, hypothesis or number of the counted statement or of its proof was altered. Layout and class: the article's own class is \texttt{amsart} with packages including amsfonts, amssymb, amsmath, amsthm, url, enumerate, inputenc, fontenc, babel, hyperref; the wrapper is the lane's pinned \texttt{amsart} editorial wrapper at 10pt on A4, which restates the theorem-like environments the retained text needs and adds the article's own macro definitions that the retained text needs (\texttt{\textbackslash FUT}, \texttt{\textbackslash LoX}, \texttt{\textbackslash Nat}, \texttt{\textbackslash RCA}, \texttt{\textbackslash W}, \texttt{\textbackslash WO}, \texttt{\textbackslash WOP}, \texttt{\textbackslash boldomega}), with extra wrapper packages (bm, bbm, dsfont, xspace, cancel, mathtools). The delivered numbering is the unit's own. Assets: no retained span contains a figure environment, a graphics-inclusion command, a PDF-page inclusion, a table or a TikZ picture; no image file is copied into this package, the package declares no asset dependency and no image file is redistributed with it. Licence: version arXiv:2401.04451v1 of this article is distributed under the Creative Commons Attribution 4.0 International licence, as recorded in the arXiv licence metadata for this exact version and shown on the abstract page https://arxiv.org/abs/2401.04451v1. A full rights notice is delivered at licenses/arxiv-2401.04451-CC-BY-4.0.txt. Characters: every retained excerpt is pure ASCII, so the delivered engine, XeLaTeX with its default Latin Modern text font, reports no missing character.
\begin{lemma}\label{lem:decribile-terms} The following is provable in $\RCA_0$: If $\alpha=(\alpha_i)_{i\in\Nat}$ is an infinite descending sequence in $\boldomega^\LoX$, then
$$\forall n\ \exists n'\ \exists m < lh(\alpha_n) \ \big(n'>n \ \land e_m(\alpha_n) >_{\LoX} e_m(\alpha_{n'})\big).$$
\end{lemma}

\begin{theorem}\label{prop:RTomega}
Let $n \geq 3, k \geq 2$. Over $\RCA_0$, $\FUT^{=n}_k$ implies $\WOP(\mathcal{X} \to \boldomega^\mathcal{X})$. Moreover, 
$$\WOP(\mathcal{X} \to \boldomega^\LoX) \leq_\W \FUT^{=n}_k.$$
\end{theorem}

\begin{proof}
Assume by way of contradiction $\neg\WO(\boldomega^\LoX)$, and let $\alpha= (\alpha_n)_{n\in \Nat}$ 
be an infinite descending sequence in $\boldomega^\LoX$. 
For this proof, it is convenient to use an $\alpha$-computable sequence $\beta$ of all the components of the terms $\alpha_n$, enumerated in order of ``appearance'', i.e. $\beta = \langle e_0(\alpha_0), e_1(\alpha_0), \dots, e_{lh(\alpha_0)-1}(\alpha_0), e_0(\alpha_1), e_1(\alpha_1), \dots \rangle$. Formally we construct such sequence by first defining $\theta : \Nat\times\Nat\to \Nat$ as follows: $\theta(n,m)= m + \sum_{k < n} lh(\alpha_k)$. The partial function $\theta$ is clearly bijective (meaning that it is a bijection between its domain of definition and its codomain). We accordingly fix functions $t: \Nat \to \Nat$ and $p:\Nat\to\Nat$ such that for each $n\in\Nat$ we have
$\theta(t(n),p(n))=n$. The sequence $(\beta_h)_{h\in\Nat}$ of components of terms in $\alpha$ is then defined
by setting $\beta_h = \alpha_{t(h),p(h)}$. Intuitively, $t(h)$ is the element of $\alpha$ from which $\beta_h$ has been ``extracted'', while $p(h)$ is the position of the component $\beta_h$ within $\alpha_{t(h)}$.

%We call $i$ \emph{decreasible} if there exists $j>i$ such that $p(j)=p(i)$ and $\beta_j > \beta_i$. In that case, we say that $j$ \emph{decreases} $i$ and, if $j$ is minimal\footnote{Il reviewer ci aveva corretto "minimal \emph{such}", correggerlo in "$j$ is the least value that decreases $i$"}, we say that $j$ \emph{min-decreases} $i$. 
%We say that $j$ \emph{min-decreases} $i<j$ if $j$ is the minimal such that $p(j)=p(i)$ and $\beta_j > \beta_i$. 
%Informally, $j$ min-decreases $i$ if $i$ is the $\beta$-index of some $e_m(\alpha_n)$ and $j$ is the $\beta$-index of the first $e_m(\alpha_{n'})$ such that $n'>n$ and $e_m(\alpha_n) > e_m(\alpha_{n'})$, i.e. $e_m(\alpha_{n'})$ is the least witness of the decreasing of $e_m(\alpha_n)$.
%We call $i$ \emph{min-decreasible} if there exists $j$ that min-decreases $i$.

We call $i$ \emph{decreasible} if there exists $j>i$ such that $p(j)=p(i)$ and $\beta_j > \beta_i$. In that case, we say that $j$ \emph{decreases} $i$ and that $j$ is a \emph{decreaser} of $i$. 

Using this terminology, Lemma \ref{lem:decribile-terms} states that each element of $\alpha$ contains at least one decreasible component.

Also, we define a number $j\in [0,r]$ \textit{important in} $S=\{n_0, \ldots,n_r\}$ if the following condition holds:
$$
(\exists i< n_0)\ \,\exists i' \in [n_{j-1}, n_j)\, \text{ s.t. } i' \text{ decreases } i \text{  and } \neg\exists i'' < n_{j-1}\, \text{  s.t. } i'' \text{  decreases } i,$$
where we set  $n_{-1}=0$.
\medskip

Now suppose that $f:\Nat \to \Nat$ is a function with the following property:

\bigskip
{\em Property P}: For all $i\in \Nat$,
if $i$ is decreasible, then it is decreased by some $j \leq f(i)$. 

\bigskip
We first show that given such an $f$ we can compute (in $f$ and $\beta$) 
an infinite descending sequence $(\sigma_i)_{i\in\Nat}$ in $\mathcal{X}$ as follows. 

\medskip
{\bf Step $0$}.  
Let $i_0$ be the least decreasible index of $\beta$, and let $j_0$ be the least decreaser of $i_0$. By Lemma \ref{lem:decribile-terms}, $\beta_{i_0}$ must be a component of $\alpha_0$, i.e. $t(i_0) = 0$, so we can find $i_0$ by just taking the least decreasible $i^* < lh(\alpha_0)$. Notice that we can decide whether $i^*$ is decreasible by just inspecting $\beta$ up to the index $f(i^*)$, since $f$ has the Property $P$.

We set $\sigma_0 = \beta_{j_0}$ and observe that $p(i) \geq p(i_0)$ for each decreasible $i>i_0$. 
Suppose otherwise as witnessed by $i^*$, and let $j^*$ be the least decreaser of $i^*$. By definition of decreaser, $p(j^*) = p(i^*) < p(i_0)$ and $\beta_{i^*} > \beta_{j^*}$. However $i_0$ is the least decreasible $\beta$-index of some component of $\alpha_0$, then $\beta_{i^*} = \beta_z$, where $z = \theta(0,p(i^*))$. Hence, $\beta_z > \beta_{j^*}$ and $p(z)=p(i^*)=p(j^*)$, but in that case $z$ would be decreasible (by $j^*$) and $p(z) < p(i_0)$, which implies $z < i_0$ since $t(z)=t(i_0)=0$, thus contradicting the minimality of $i_0$.

\bigskip
{\bf Step $s+1$}. Suppose $i_s, j_s, \sigma_s$ are defined, $(\sigma_t)_{t \leq s}$ is decreasing in $\mathcal{X}$ and 
$p(i) \geq p(i_s)$ for each decreasible $i > i_s$. 

Let $i_{s+1}$ be the least decreasible index of $\beta$ larger than or equal to $j_s$, and let $j_{s+1}$ be the least decreaser of $i_{s+1}$. By Lemma \ref{lem:decribile-terms} and the fact that no decreasible $i > i_s$ can have $p(i) < p(i_s)=p(j_s)$, $\beta_{i_{s+1}}$ must be a component of $\alpha_{t(j_s)}$, namely the leftmost whose $\beta$-index is decreasible. So we can find $i_{s+1}$ 
by just taking the least decreasible $i^* \in [j_s,\ j_s + lh(\alpha_{t(j_s)}) - p(j_s))$. Notice that we can decide whether $i^*$ is decreasible by just inspecting $\beta$ up to the index $f(i^*)$, since $f$ has the Property $P$.

We set $\sigma_{s+1} = \beta_{j_{s+1}}$. As we noted above, $\beta_{i_{s+1}}$ must be either $\beta_{j_s}$ or a component of $\alpha_{t(j_s)}$ on the right of $\beta_{j_s}$, i.e. $t(i_{s+1}) = t(j_s)$ and $p(i_{s+1}) \geq p(j_s)$, so $\beta_{j_s} \geq \beta_{i_{s+1}}$. Then, $\sigma_s > \sigma_{s+1}$ because $\sigma_s = \beta_{j_s} \geq \beta_{i_{s+1}} > \beta_{j_{s+1}} = \sigma_{s+1}$. Finally, we observe that the last part of the inductive invariant is guaranteed as well, since $p(i) \geq p(i_{s+1})$ for each decreasible $i>i_{s+1}$. 
Suppose otherwise as witnessed by $i^*$, and let $j^*$ be the least decreaser of $i^*$. By definition of decreaser, $p(j^*) = p(i^*) < p(i_{s+1})$ and $\beta_{i^*} > \beta_{j^*}$. However $\beta_{i_{s+1}}$ is the leftmost component of $\alpha_{t(j_s)}$ whose $\beta$-index is decreasible, then $\beta_{i^*} = \beta_z$, where $z = \theta({t(j_s)},p(i^*))$. Hence, $\beta_z > \beta_{j^*}$, but in that case $z$ would be decreasible (by $j^*$) and $p(z) = p(i^*) < p(i_{s+1})$, which implies $z < i_{s+1}$ since $t(z)=t(j_s)=t(i_{s+1})$, thus contradicting the minimality of $i_{s+1}$.


\bigskip
We now show how to obtain a function satisfying the Property $P$ from a solution of $\FUT^{=n}_k$ for
a suitable colouring. Let $g: FIN(\Nat^+) \rightarrow k$ as follows:
$$
g(S) = card\{j\ |\ j \text{ is important in }  S\} \text{ mod } k.
$$

\medskip

By $\FUT^{=n}_k$ let $\mathcal{B} = \{B_0 < B_1 < B_2 < \ldots\}$ be an infinite block sequence such that $FU^{=n}(\mathcal{B})$ is monochromatic under $g$, and let $c<k$ be the colour of $\mathcal{B}$.

\begin{claim} \label{clm:shorter-union}
Given $S_0\!<\!\ldots\!<\!S_{n-3}$ in $\mathcal{B}$, there exists $S \in \mathcal{B}$ such that $S_{n-3} < S$ and $g(S_0 \cup \ldots \cup S_{n-3} \cup S) = c$.
\end{claim}

Fix $S_0<\ldots<S_{n-3}$ in $\mathcal{B}$ and let $\ell$ be the actual upper bound of the minimal indexes decreasing all the decreasible $j < \min(S_0)$. By this we mean the $\ell$ given by the following instance of strong $\Sigma^0_1$-bounding (in $\RCA_0$): 
$$
\forall m \,\exists \ell \,\forall j < m \,( \exists d \,(d \text{ decreases } j) \to \exists d < \ell \,(d \text{ decreases } j))),
$$
where we can take $m=\min(S_0)$. 

Since $\mathcal{B}$ is an infinite block sequence, there exists $S \in \mathcal{B}$ such that $\min(S) > \ell > \max(S_{n-3})$. Then, for any $T \in \mathcal{B}$ with $S < T$, we have that $g(S_0 \cup \ldots \cup S_{n-3} \cup S) = g(S_0 \cup \ldots \cup S_{n-3} \cup S \cup T)$, since no elements in $T$ are important in $S_0 \cup \ldots \cup S_{n-3} \cup S \cup T$. Also, by monochromaticity of $\mathcal{B}$, $g(S_0 \cup \ldots \cup S_{n-3} \cup S \cup T) = c$, so $g(S_0 \cup \ldots \cup S_{n-3} \cup S) = c$, hence proving the Claim.

\bigskip

Now, we define $f: \Nat \to \Nat$ as $f(i) = \max(B_q)$, where $q$ is minimal such that $g(B_p \cup B_{p+1} \cup B_{p+2} \cup \ldots \cup B_{p+n-3} \cup B_q) =c$, with $p$ minimal such that $i < \min(B_p)$ and $B_p < B_{p+1} < B_{p+2} < \ldots < B_{p+n-3} < B_q$. Notice that $q$ exists by Claim \ref{clm:shorter-union}. Also, $f$ has the Property $P$, i.e., each decreasible $i$ is decreased by some $j \leq f(i)$.

In order to prove this, assume by way of contradiction that $i$ is decreasible and $p,q$ are minimal such that $i < \min(B_p)$, $B_p < B_{p+1} < \ldots < B_{p+n-3} < B_q$ and $g(B_p \cup B_{p+1} \cup \ldots \cup B_{p+n-3} \cup B_q) = c$, but $i$ is decreasible only by numbers larger than $f(i) = \max(B_q)$. 

By strong $\Sigma^0_1$-bounding, let $\ell$ be the actual upper bound of the minimal indexes decreasing all the decreasible $j < \min(B_p)$. Since $\mathcal{B}$ is an infinite block sequence, there exists $B \in \mathcal{B}$ such that $\min(B) > \ell > \max(B_q)$. Now, consider:
$$
B_p \cup \ldots \cup B_{p+n-3} \cup B_q \cup B = \{b_0, \ldots, b_r, b_{r+1}, \ldots b_t\},$$
where $b_0 = \min(B_p) = \min(B_p \cup \ldots \cup B_{p+n-3} \cup B_q \cup B)$, $b_r = \max(B_q)$ and $b_{r+1} = \min(B)$. Clearly, $j \leq t$ is important in $B_p \cup \ldots \cup B_{p+n-3} \cup B_q \cup B$ if and only if either $j \leq r$ and $j$ is important in $B_p \cup \ldots \cup B_{p+n-3} \cup B_q$, or $j = r + 1$; hence, $g(B_p \cup \ldots \cup B_{p+n-3} \cup B_q) \neq g(B_p \cup \ldots \cup B_{p+n-3} \cup B_q \cup B) = c$, contra our assumption that $g(B_p \cup \ldots \cup B_{p+n-3} \cup B_q) = c$.
\end{proof}

\end{document}
